\chapter{信号与波}\label{ch:g10:signals-and-waves}

你曾收到的每一条消息，都以\omterm{def:g10:signals-and-waves:signal}{信号}的形式抵达：一段\omterm{def:g10:signals-and-waves:sound}{声音}压在你的鼓膜上，
这一页文字则以无线电闪烁的形式传到你面前。背后
千变万化之下，只有一个想法——随时间变化的\omterm{def:g10:orders-of-magnitude:unit}{物理量}——以及一套
工具：\omterm{def:g10:signals-and-waves:period}{周期}、\omterm{def:g10:signals-and-waves:frequency}{频率}、\omterm{def:g10:signals-and-waves:amplitude}{振幅}，以及测定延迟的技巧。
凭借这些，船只可测定海底，医生可数出心跳。

\section{信号携带信息}

\begin{definition}[信号]\label{def:g10:signals-and-waves:signal}
\emph{信号}\index{信号}是随时间变化、用以携带信息的\omterm{def:g10:orders-of-magnitude:unit}{物理量}：
耳膜处的气压、话筒接线柱两端的电压、
\omterm{def:g10:refraction:fiber}{光纤}中光的亮度。
研究信号时，应将其记录为时间的函数，并从
曲线上读取信息。
\end{definition}

一通电话中，同一信息会多次更换载体：
先是空气中的\omterm{def:g10:pressure:pressure}{压强}，再是话筒中的电压，然后是
天线处的无线电波——对端则反向进行。一套词汇
即可统一描述这些形式；本章即构建它。

\section{周期信号}

\begin{definition}[周期信号、周期]\label{def:g10:signals-and-waves:period}
若\omterm{def:g10:signals-and-waves:signal}{信号}在规则的时间间隔上完全重复，则称其为\emph{周期性的}\index{周期信号}。
\emph{周期}\index{周期} $T$ 是
一个完整循环所经历的时长，以秒为\omterm{def:g10:orders-of-magnitude:unit}{单位}。
\end{definition}

\begin{definition}[频率]\label{def:g10:signals-and-waves:frequency}
\omterm{def:g10:signals-and-waves:period}{周期信号}的\emph{频率}\index{频率}是
每秒的循环次数：
\[
f = \frac{1}{T} .
\]
其\omterm{def:g10:orders-of-magnitude:unit}{单位}为每秒一个循环，即\emph{赫兹}\index{赫兹}
（\unit{Hz}），常用倍数为 \unit{kHz}、\unit{MHz}、\unit{GHz}。
\end{definition}

\begin{definition}[振幅]\label{def:g10:signals-and-waves:amplitude}
\omterm{def:g10:signals-and-waves:period}{周期信号}的\emph{振幅}\index{振幅}是其
相对静止值的最大偏差。对于围绕零值对称振荡的
电压，振幅即峰值；从波峰到波谷的量度为
\emph{两倍}振幅（即峰峰值）。
\end{definition}

\begin{example}[市电与音乐会标准音]\label{ex:g10:signals-and-waves:mains}
市电电压以 $f = \qty{50}{Hz}$ 振荡，其\omterm{def:g10:signals-and-waves:period}{周期}为
$T = 1/f = \qty{0.020}{s} = \qty{20}{ms}$。标准音 A 是一种
\omterm{def:g10:signals-and-waves:frequency}{频率}为 \qty{440}{Hz} 的\omterm{def:g10:pressure:pressure}{压强}\omterm{def:g10:signals-and-waves:signal}{信号}：一个\omterm{def:g10:signals-and-waves:period}{周期}持续
$T = 1/440 = \qty{2.3}{ms}$。\omterm{def:g10:signals-and-waves:frequency}{频率}越高，\omterm{def:g10:signals-and-waves:period}{周期}越短：$f$ 与 $T$
互为倒数。
\end{example}

\begin{omfigure}
\begin{tikzpicture}[font=\small]
  \draw[-Stealth] (-0.3,0) -- (7.7,0) node[below] {$t$};
  \draw[-Stealth] (0,-1.7) -- (0,2.0) node[left] {信号};
  \draw[omDef, thick, smooth, samples=120, domain=0:7.3]
    plot (\x, {1.3*sin(120*\x)});
  \draw[dashed, thin] (0.75,1.3) -- (0.75,1.6) (3.75,1.3) -- (3.75,1.6);
  \draw[omThm, Stealth-Stealth] (0.75,1.6) -- (3.75,1.6)
    node[midway, above] {周期 $T$};
  \draw[omProp, Stealth-Stealth] (5.25,0) -- (5.25,-1.3)
    node[midway, right] {振幅};
\end{tikzpicture}

{\small \omterm{def:g10:signals-and-waves:period}{周期信号}：图样每隔 $T$ 秒重复一次，而
\omterm{def:g10:signals-and-waves:amplitude}{振幅}是相对静止值的最大偏差。}
\end{omfigure}

\section{读取示波图}

\begin{definition}[示波器]\label{def:g10:signals-and-waves:oscilloscope}
\emph{示波器}\index{示波器}描绘电压随时间的变化；
其带刻度屏幕上的曲线称为\emph{示波图}%
\index{示波图}。两个设置可把屏幕格数换算为\omterm{def:g10:orders-of-magnitude:unit}{物理量}：
其一是\emph{时基}\index{时基}（每水平格的秒数），
另一是竖直方向的\emph{增益}\index{增益}（每竖直格
的伏特数）。
\end{definition}

\begin{method}[从屏幕读数到数值]\label{met:g10:signals-and-waves:oscillogram}
\begin{enumerate}
  \item 数出一个完整\omterm{def:g10:signals-and-waves:period}{周期}所跨的水平格数
        （以波峰到波峰最稳妥）；乘以\omterm{def:g10:signals-and-waves:oscilloscope}{时基}得到
        $T$，再求 $f = 1/T$。
  \item 数出从中线到某一波峰的竖直格数，
        乘以\omterm{def:g10:signals-and-waves:oscilloscope}{增益}即得\omterm{def:g10:signals-and-waves:amplitude}{振幅}（波峰到波谷给出峰峰值，
        须取其一半）。
\end{enumerate}
\end{method}

\begin{example}[完整读数示例]\label{ex:g10:signals-and-waves:reading}
下图\omterm{def:g10:signals-and-waves:oscilloscope}{示波图}中，\omterm{def:g10:signals-and-waves:oscilloscope}{时基}为每格 \qty{5}{ms}，
\omterm{def:g10:signals-and-waves:oscilloscope}{增益}为每格 \qty{2}{V}。一个完整\omterm{def:g10:signals-and-waves:period}{周期}跨越 $4.0$ 格：
$T = 4.0 \times \qty{5}{ms} = \qty{20}{ms}$，且
$f = 1/\qty{0.020}{s} = \qty{50}{Hz}$——正是市电\omterm{def:g10:signals-and-waves:frequency}{频率}。一个波峰位于
坐标轴上方 $3.0$ 格处：\omterm{def:g10:signals-and-waves:amplitude}{振幅}为
$3.0 \times \qty{2}{V} = \qty{6.0}{V}$。
\end{example}

\begin{omfigure}
\begin{tikzpicture}[x=0.62cm, y=0.62cm, font=\small]
  \draw[gray!35, thin] (0,0) grid (10,8);
  \draw[gray!80] (0,4) -- (10,4);
  \draw[gray!80] (5,0) -- (5,8);
  \draw[omDef, thick, smooth, samples=120, domain=0:10]
    plot (\x, {4 + 3*sin(90*\x)});
  \draw[dashed, thin] (1,7) -- (1,7.6) (5,7) -- (5,7.6);
  \draw[omThm, Stealth-Stealth] (1,7.6) -- (5,7.6)
    node[midway, above] {$4.0$ 格};
  \draw[omProp, Stealth-Stealth] (7,4) -- (7,1)
    node[midway, right] {$3.0$ 格};
\end{tikzpicture}

{\small \omterm{def:g10:signals-and-waves:oscilloscope}{示波图}，\omterm{def:g10:signals-and-waves:oscilloscope}{时基} \qty{5}{ms}/格，\omterm{def:g10:signals-and-waves:oscilloscope}{增益} \qty{2}{V}/格：
一个\omterm{def:g10:signals-and-waves:period}{周期}占 $4.0$ 格（$T = \qty{20}{ms}$，$f = \qty{50}{Hz}$），
波峰高出 $3.0$ 格（\omterm{def:g10:signals-and-waves:amplitude}{振幅} \qty{6.0}{V}）。}
\end{omfigure}

\section{声音与电磁信号}

\begin{definition}[声音]\label{def:g10:signals-and-waves:sound}
\emph{声音}\index{声音}是一种\omterm{def:g10:pressure:pressure}{压强}\omterm{def:g10:signals-and-waves:signal}{信号}：振动物体——
琴弦、薄膜、声带——有节奏地推动空气，
密部以约 \qty{340}{m/s} 的速率向外传播。麦克风
把到达的\omterm{def:g10:pressure:pressure}{压强}转换为电压；耳膜则把它转换为
神经冲动。声音需要介质：在真空钟罩下敲响的铃铛，
随着空气被抽出而逐渐无声。
\end{definition}

\begin{definition}[波]\label{def:g10:signals-and-waves:wave}
行进的扰动——\omterm{def:g10:signals-and-waves:sound}{声音}的密部图样、池塘上的涟漪、
无线电脉冲——都是一种\emph{波}\index{波}。若有介质，
介质本身保持不动：每一小块水面原地起伏，
而涟漪掠过池塘。波的几何特征——\omterm{def:g10:light-spectra:wavelength}{波长}、
干涉、衍射——将在后面一章讨论；本章
只需用到它们的速率。
\end{definition}

\begin{definition}[超声波与次声波]\label{def:g10:signals-and-waves:ultrasound}
人耳能听到约 \qty{20}{Hz} 至 \qty{20}{kHz} 之间的\omterm{def:g10:signals-and-waves:sound}{声音}
（上限随年龄增长而降低）。高于 \qty{20}{kHz} 的\omterm{def:g10:signals-and-waves:sound}{声音}是
\emph{超声波}\index{超声波}——蝙蝠捕食时使用 \qty{50}{kHz} 的声波，
医用扫描仪以几兆赫的\omterm{def:g10:signals-and-waves:frequency}{频率}为人体成像。低于
\qty{20}{Hz} 的\omterm{def:g10:signals-and-waves:sound}{声音}是\emph{次声波}\index{次声波}，能被大象
和地震仪感知。
\end{definition}

\begin{proposition}[电磁信号]\label{prop:g10:signals-and-waves:em}
无线电波、Wi-Fi 与光都是\emph{电磁
信号}\index{电磁信号}：同属一族，均可在真空中传播，
无需介质——且速率相同，都等于
在\cref{ch:g10:refraction}中所述的光速，
\[
c = \qty{3.00e8}{m/s} \quad \text{（在空气中几乎相同）。}
\]
\end{proposition}
\admitted

\begin{remark}\label{rem:g10:signals-and-waves:honest}
这一点是一项实验事实——阳光穿过 $150$ 百万
千米的空旷空间；此事实将在本科二年级分册中依电磁学定律予以推导。
请记住这一对比：\omterm{def:g10:signals-and-waves:sound}{声音}是介质
\emph{的}振动，没有介质便会消失；\omterm{prop:g10:signals-and-waves:em}{电磁信号}则携带自身的
场穿越真空，速率快一百万倍。
\end{remark}

\begin{omfigure}
\begin{tikzpicture}[x=1.3cm, font=\small]
  \fill[omDef!15] (1.3,-0.14) rectangle (4.3,0.14);
  \draw[-Stealth] (-0.3,0) -- (7.7,0);
  \foreach \x/\l in {0/{$1$}, 1/{$10$}, 2/{$10^2$}, 3/{$10^3$},
      4/{$10^4$}, 5/{$10^5$}, 6/{$10^6$}, 7/{$10^7$}}
    \draw (\x,0.07) -- (\x,-0.07) node[below] {\l};
  \node[below] at (7.55,-0.07) {\unit{Hz}};
  \node[gray] at (0.4,-0.75) {次声波};
  \node[omDef] at (2.8,-0.75) {可听声};
  \node[gray] at (5.8,-0.75) {超声波};
  \draw[omThm] (1.7,0.14) -- (1.7,0.4) node[above] {市电};
  \draw[omThm] (2.64,0.14) -- (2.64,0.95) node[above] {标准音 A};
  \draw[omThm] (4.7,0.14) -- (4.7,0.4) node[above] {蝙蝠};
  \draw[omThm] (6.7,0.14) -- (6.7,0.95) node[above] {医用扫描};
\end{tikzpicture}

{\small \omterm{def:g10:signals-and-waves:sound}{声音}\omterm{def:g10:signals-and-waves:frequency}{频率}采用对数刻度（每一步 $\times 10$）。
人耳可听范围为 \qty{20}{Hz} 至 \qty{20}{kHz}；其下为\omterm{def:g10:signals-and-waves:ultrasound}{次声波}，
其上为\omterm{def:g10:signals-and-waves:ultrasound}{超声波}。}
\end{omfigure}

\section{回声：由延迟求距离}

以速率 $v$ 传播的\omterm{def:g10:signals-and-waves:signal}{信号}，走过的距离为 $d = v\,t$，所用时间为
$t$：每一段延迟都是伪装的距离。更好的做法是：向一个障碍物发出短脉冲，
并计时其\emph{回声}\index{回声}——往返路程为 $2d$。

\begin{method}[回声测距]\label{met:g10:signals-and-waves:echo}
\begin{enumerate}
  \item 发出一个短脉冲；测量\omterm{rem:g10:signals-and-waves:honest}{回声}的往返时间 $t$。
  \item 取\omterm{def:g10:signals-and-waves:signal}{信号}\emph{在所经过的介质中}的速率 $v$。
  \item 脉冲经过的路程为 $2d = v\,t$，故 $d = \dfrac{v\,t}{2}$。
\end{enumerate}
水中的声脉冲用于\emph{声呐}\index{声呐}
（$v \approx \qty{1500}{m/s}$）；空气中的无线电脉冲用于
\emph{雷达}\index{雷达}（$v = c$）；体内的超声脉冲用于
医用扫描仪（软组织中 $v \approx \qty{1540}{m/s}$）。
\end{method}

\begin{example}[声呐]\label{ex:g10:signals-and-waves:sonar}
一艘船的\omterm{met:g10:signals-and-waves:echo}{声呐}发出脉冲，海底\omterm{rem:g10:signals-and-waves:honest}{回声}
经过 $t = \qty{0.80}{s}$ 后返回：水深为
$d = \frac{1500 \times 0.80}{2} = \qty{600}{m}$。忘记因数
$2$ 会使算得的海洋深度翻倍。
\end{example}

\begin{omfigure}
\begin{tikzpicture}[font=\small]
  \fill[blue!8] (-3.5,0) -- (3.5,0) -- (3.5,-2.95) -- (2.6,-3.1)
    -- (1.4,-2.9) -- (0.2,-3.05) -- (-1.2,-2.85) -- (-2.3,-3.1)
    -- (-3.5,-2.95) -- cycle;
  \draw[thick, blue!60] (-3.5,0) -- (3.5,0);
  \draw[thick] (-3.5,-2.95) -- (-2.3,-3.1) -- (-1.2,-2.85)
    -- (0.2,-3.05) -- (1.4,-2.9) -- (2.6,-3.1) -- (3.5,-2.95);
  \fill[pattern=north east lines] (-3.5,-2.95) -- (-2.3,-3.1)
    -- (-1.2,-2.85) -- (0.2,-3.05) -- (1.4,-2.9) -- (2.6,-3.1)
    -- (3.5,-2.95) -- (3.5,-3.4) -- (-3.5,-3.4) -- cycle;
  \fill[gray!70] (-0.9,0.3) -- (0.9,0.3) -- (0.55,-0.25)
    -- (-0.55,-0.25) -- cycle;
  \fill[gray!50] (-0.25,0.3) rectangle (0.25,0.62);
  \draw[-Stealth, omDef, thick] (-0.35,-0.35) -- (-0.35,-2.75)
    node[midway, left] {发出，$t/2$};
  \draw[-Stealth, omThm, thick] (0.35,-2.75) -- (0.35,-0.35)
    node[midway, right] {回声，$t/2$};
  \draw[Stealth-Stealth] (2.9,-0.05) -- (2.9,-2.9)
    node[midway, right] {$d = \dfrac{v\,t}{2}$};
\end{tikzpicture}

{\small \omterm{met:g10:signals-and-waves:echo}{声呐}脉冲往返传播：测得的延迟 $t$ 对应
$2d$ 的路程，因此深度 $d = v\,t/2$，其中 $v \approx \qty{1500}{m/s}$
为海水中的速率。}
\end{omfigure}

\begin{example}[雷达]\label{ex:g10:signals-and-waves:radar}
机场\omterm{met:g10:signals-and-waves:echo}{雷达}在发出脉冲后
$\qty{0.30}{ms}$ 收到飞机的回波：
$d = \frac{\qty{3.00e8}{m/s} \times \qty{3.0e-4}{s}}{2} = \qty{45}{km}$.
因为 $c$ 极大，\omterm{met:g10:signals-and-waves:echo}{雷达}所用的时钟精度要达到微秒量级，
对光的计时是精密工作。
\end{example}

\begin{example}[心脏作为信号]\label{ex:g10:signals-and-waves:ecg}
每次心跳都会使胸部产生一个微小的电压脉冲；电极
将其记录为\emph{心电图}\index{心电图}（ECG），一种近似\omterm{def:g10:signals-and-waves:period}{周期性的}
\omterm{def:g10:signals-and-waves:signal}{信号}。若高峰之间相隔 $T = \qty{0.75}{s}$，则心脏
跳动\omterm{def:g10:signals-and-waves:frequency}{频率}为 $f = 1/0.75 = \qty{1.33}{Hz}$，即 $1.33 \times 60 = 80$
次/分钟。从图线上读取 $T$，
正是心脏病专家测量脉搏的方法。
\end{example}

\section{练习}

\begin{exercise}[$\star$]\label{exo:g10:signals-and-waves:1}
求\omterm{def:g10:signals-and-waves:period}{周期}为下列各值的\omterm{def:g10:signals-and-waves:period}{周期信号}的\omterm{def:g10:signals-and-waves:frequency}{频率}：
(a) $T = \qty{20}{ms}$；(b) $T = \qty{4.0}{ms}$；
(c) $T = \qty{50}{\micro s}$。
\end{exercise}

\begin{exercise}[$\star$]\label{exo:g10:signals-and-waves:2}
求下列各项的\omterm{def:g10:signals-and-waves:period}{周期}：(a) \qty{50}{Hz} 的市电；(b) \qty{440}{Hz}
的标准 A 音；(c) \qty{2.4}{GHz} 的 Wi-Fi \omterm{def:g10:signals-and-waves:signal}{信号}。
\end{exercise}

\begin{exercise}[$\star$]\label{exo:g10:signals-and-waves:3}
判断下列声波属于\omterm{def:g10:signals-and-waves:ultrasound}{次声波}、可闻声还是\omterm{def:g10:signals-and-waves:ultrasound}{超声波}：\qty{12}{Hz}；
\qty{440}{Hz}；\qty{18}{kHz}；\qty{40}{kHz}（汽车泊车传感器）；
\qty{5}{MHz}（医用探头）。
\end{exercise}

\begin{exercise}[$\star$]\label{exo:g10:signals-and-waves:4}
在\omterm{def:g10:signals-and-waves:oscilloscope}{示波器}上，\omterm{def:g10:signals-and-waves:oscilloscope}{时基}设为每格 \qty{2}{ms}，电压灵敏度设为每格 \qty{0.5}{V}，
一个\omterm{def:g10:signals-and-waves:period}{周期}占 $5.0$ 格，波峰位于
坐标轴上方 $4.0$ 格。求 $T$、$f$ 和\omterm{def:g10:signals-and-waves:amplitude}{振幅}。
\end{exercise}

\begin{exercise}[$\star$]\label{exo:g10:signals-and-waves:5}
先看到闪电，\qty{6.0}{s} 后才听到雷声。
闪电击中的地点有多远（$v_{\text{声音}} = \qty{340}{m/s}$）？为何可以忽略
光的传播时间？
\end{exercise}

\begin{exercise}[$\star\star$]\label{exo:g10:signals-and-waves:6}
\omterm{met:g10:signals-and-waves:echo}{声呐}回波在 \qty{0.60}{s} 后从海底返回
（海水中 $v = \qty{1500}{m/s}$）。为什么乘积 $v\,t$ 必须
减半？求深度。若是在一个
\qty{1200}{m} 深的海沟上，回波需要多长时间？
\end{exercise}

\begin{exercise}[$\star\star$]\label{exo:g10:signals-and-waves:7}
机场\omterm{met:g10:signals-and-waves:echo}{雷达}在 \qty{0.24}{ms} 后收到一个飞机回波，
恰好 \qty{1.0}{s} 之后又经 \qty{0.238}{ms} 收到第二个回波。计算
两段距离，再求飞机的速率。它是否正在靠近？
\end{exercise}

\begin{exercise}[$\star\star$]\label{exo:g10:signals-and-waves:8}
\omterm{ex:g10:signals-and-waves:ecg}{心电图}以 \qty{25}{mm/s} 的速度打印，高峰之间相隔
\qty{20}{mm}。求\omterm{def:g10:signals-and-waves:period}{周期}，以及以次/分钟计的心率。
心动过速指每分钟跳动超过 $100$ 次。在这种纸上，
高峰间距低于多少时会表现为心动过速？
\end{exercise}

\begin{exercise}[$\star\star$]\label{exo:g10:signals-and-waves:9}
需在屏幕上显示\omterm{def:g10:signals-and-waves:amplitude}{振幅}为 \qty{3.0}{V}、\omterm{def:g10:signals-and-waves:frequency}{频率}为 \qty{200}{Hz} 的\omterm{def:g10:signals-and-waves:signal}{信号}，
屏幕水平方向 $10$ 格、竖直方向 $8$ 格（中线以上 $4$ 格
），选择\omterm{def:g10:signals-and-waves:oscilloscope}{时基}显示约两个完整\omterm{def:g10:signals-and-waves:period}{周期}，并选择
\omterm{def:g10:signals-and-waves:oscilloscope}{增益}，使\omterm{def:g10:signals-and-waves:signal}{信号}占据屏幕大部分且不出现削波。
\end{exercise}

\begin{exercise}[$\star\star$]\label{exo:g10:signals-and-waves:10}
医用探头向腹部发射一个超声脉冲
（软组织中 $v = \qty{1540}{m/s}$），并先后收到回波，
时间分别为 \qty{40}{\micro s}（器官前壁）和 \qty{60}{\micro s}（后
壁）。求各壁的深度以及该器官的厚度。
\end{exercise}

\begin{exercise}[$\star\star$]\label{exo:g10:signals-and-waves:11}
一场音乐会正在现场直播。谁先听到鼓点：离舞台
\qty{30}{m} 的现场听众，还是 \qty{3000}{km} 外的广播听众？
两者相差多少时间？
\end{exercise}

\begin{exercise}[$\star\star\star$]\label{exo:g10:signals-and-waves:12}
蝙蝠发射持续 \qty{3.0}{ms} 的超声脉冲，发射期间无法听到
回波。飞蛾在距离小于多少时无法被探测到？
距离 \qty{1.7}{m} 的飞蛾，回波延迟是多少？
蝙蝠在最后逼近时为何必须缩短脉冲？
\end{exercise}

\begin{exercise}[$\star\star\star$]\label{exo:g10:signals-and-waves:13}
你面对悬崖拍手，\qty{1.5}{s} 后听到\omterm{rem:g10:signals-and-waves:honest}{回声}。悬崖距离
多远？你沿直线朝悬崖走 \qty{100}{m} 后，\omterm{rem:g10:signals-and-waves:honest}{回声}延迟是多少？耳朵
在约 \qty{0.1}{s} 以内，耳朵无法区分拍手与\omterm{rem:g10:signals-and-waves:honest}{回声}。在多近的距离内，
\omterm{rem:g10:signals-and-waves:honest}{回声}会融入拍手声而消失？
\end{exercise}

\begin{exercise}[$\star\star\star$]\label{exo:g10:signals-and-waves:14}
在每格 \qty{5}{ms} 的\omterm{def:g10:signals-and-waves:oscilloscope}{时基}下，一个\omterm{def:g10:signals-and-waves:signal}{信号}\omterm{def:g10:signals-and-waves:period}{周期}占 $3.0$ 格
，屏幕共 $10$ 格。
\begin{enumerate}
  \item 求 $T$ 和 $f$。
  \item 将\omterm{def:g10:signals-and-waves:oscilloscope}{时基}改为每格 \qty{2}{ms}，此时一个\omterm{def:g10:signals-and-waves:period}{周期}
        占多少格？
  \item 在每种设置下，屏幕上能容纳多少个
        \emph{完整}\omterm{def:g10:signals-and-waves:period}{周期}？
\end{enumerate}
\end{exercise}

\begin{exercise}[$\star\star\star$]\label{exo:g10:signals-and-waves:15}
用“ping”测量一个互联网数据包往返
一台位于 \qty{1200}{km} 外的服务器的最短可能往返
时间是多少（\omterm{def:g10:signals-and-waves:signal}{信号}以 $c$ 传播）？在\omterm{def:g10:refraction:fiber}{光纤}中，光速约为
\qty{2.0e8}{m/s}，请重新计算。实测 ping 值为 \qty{40}{ms}，请给出两个
理由，说明其为何大于你的计算结果。
\end{exercise}

\section{问题：声呐、风暴与心电图}

\begin{problem}[{周末问题——读懂世界的信号：一条
定律，$d = v\,t$，用来测定风暴距离、绘制海床、计数心跳，并
探究 GPS 的工作原理}]
\label{pb:g10:signals-and-waves:1}
一艘渔船度过了一个暴风雨的夜晚：地平线上闪着电光，
\omterm{met:g10:signals-and-waves:echo}{声呐}正在扫描海底，船长的心跳正描在医生的
纸带上，GPS 接收机正在收听\omterm{def:g10:universal-gravitation:satellite}{卫星}。四种仪器，一条
定律。取空气中 $v_{\text{声}} = \qty{340}{m/s}$，海水中
\qty{1500}{m/s}，且 $c = \qty{3.00e8}{m/s}$。

\medskip
\noindent\textbf{第一部分——风暴。}
\begin{enumerate}
  \item 一次闪电之后 \qty{9.0}{s}，传来雷声。闪电发生处离船多远？
  \item 计算光在该距离上的传播时间，并
        说明将闪电视为瞬时发生的理由。
  \item 水手数出闪电与雷声之间的秒数，再除以
        三，即得千米数。请说明这一规则的依据。
  \item 五分钟后，又一次闪电的间隔为 \qty{6.0}{s}。此时距离多远？
        求风暴平均靠近的速率，分别以 \unit{m/s} 和 \unit{km/h} 表示。
  \item 雷声隆隆而非劈啪作响：闪电通道长达数千米，
        因此其各部分离船的距离各不相同。
        若闪电通道离船从 \qty{2.0}{km} 延伸至 \qty{5.0}{km}，
        雷声会持续多久？
\end{enumerate}

\medskip
\noindent\textbf{第二部分——\omterm{met:g10:signals-and-waves:echo}{声呐}。}
\begin{enumerate}[resume]
  \item \omterm{met:g10:signals-and-waves:echo}{声呐}发出脉冲，海床的回波在 \qty{0.90}{s} 后返回。水深多少？
  \item 在大陆架上，回波时间缩短为 \qty{0.20}{s}。水深多少？
  \item 一次脉冲返回\emph{两}个回波，时间分别为 \qty{0.16}{s} 和
        \qty{0.20}{s}。请解释这两个回波的含义，并说明鱼群
        在离海底多高处游动。
  \item \omterm{met:g10:signals-and-waves:echo}{声呐}每隔 \qty{0.50}{s} 发射一次脉冲，能无混淆地测量的最大深度为多少，
        超过该深度又会出现什么问题？
  \item 若在空气中，同样的 \qty{0.60}{s} 延迟对应多远的距离？
        寓意：把延迟换算成距离之前，必须先知道
        什么？
\end{enumerate}

\medskip
\noindent\textbf{第三部分——\omterm{ex:g10:signals-and-waves:ecg}{心电图}。}
\begin{enumerate}[resume]
  \item 船长的\omterm{ex:g10:signals-and-waves:ecg}{心电图}波峰相隔 \qty{0.86}{s}。\omterm{def:g10:signals-and-waves:frequency}{频率}是多少？
        每分钟跳动多少次？
  \item 纸带以 \qty{25}{mm/s} 的速率前进，医生
        测得的波峰间距为多少？
  \item 收网之后，间距变为 \qty{15}{mm}。新的心率是多少？
  \item 一名年轻的甲板水手是训练有素的赛艇运动员，静息时
        心率为每分钟 $50$ 次。\omterm{def:g10:signals-and-waves:period}{周期}是多少？纸带上的间距呢？
  \item \omterm{ex:g10:signals-and-waves:ecg}{心电图}是否严格\omterm{def:g10:signals-and-waves:period}{周期}性？那么，医生从
        描迹上读取的是什么？
\end{enumerate}

\medskip
\noindent\textbf{第四部分——借光返航。}
\begin{enumerate}[resume]
  \item 渔船用无线电联系距离 \qty{50}{km} 的港口，
        信息需要多长时间到达？若改用\omterm{def:g10:signals-and-waves:sound}{声音}，又需多长时间？
  \item 一颗 GPS \omterm{def:g10:universal-gravitation:satellite}{卫星}在约 \qty{20200}{km} 的上空运行，其\omterm{def:g10:signals-and-waves:signal}{信号}
        需要多长时间到达渔船？
  \item GPS 把时间转换为位置。
        一个 \qty{1.0}{\micro s} 的时钟误差会造成多大的距离误差？精度为
        \qty{3.0}{m} 的定位，需要多高的时间精度？
  \item 即使原则上，为什么也不可能存在基于\omterm{def:g10:signals-and-waves:sound}{声音}的 GPS？
  \item 总结：列出今晚用到的三种速率，以及四种仪器背后
        共同的那一条定律，然后用一句话
        概括测距的习惯。
\end{enumerate}
\end{problem}
